\subsection{曲線を直線に置き換える}

定積分を近似する最初の方法は，区間を細かくして長方形の面積を足すことであった。両端の関数値を使えば，長方形の代わりに台形を使える。区間$[a,b]$で，両端の点を通る直線は
\[
 L(x)=\frac{b-x}{b-a}f(a)+\frac{x-a}{b-a}f(b)
\]
である。この直線の下の符号付き面積は
\[
 T=\int_a^bL(x)\dd x=\frac{b-a}{2}\{f(a)+f(b)\}
\]
となる。これが一つの区間に対する台形公式である。

図\ref{図:台形近似}では，曲線が下に凸なので弦が曲線より上にある。この場合，台形公式の値は積分の値より大きくなる。ただし，関数によって凸凹が変わる場合もある。そこで，図だけでなく二階微分による誤差の式を求めよう。以下の台形公式の誤差評価では，$f$ は $[a,b]$ で二回連続微分可能であるとする。

\begin{figure}[htbp]
\centering
\begin{tikzpicture}[x=1.7cm,y=1.4cm]
 \fill[blue!10] (0.5,0)--(2.5,0)--(2.5,1.9625)--(0.5,0.4625)--cycle;
 \draw[-{Latex}] (-0.1,0)--(2.9,0) node[right] {$x$};
 \draw[-{Latex}] (0,-0.1)--(0,2.35) node[above] {$y$};
 \draw[dashed] (0.5,0)--(0.5,0.4625);
 \draw[dashed] (2.5,0)--(2.5,1.9625);
 \draw[thick,orange!80!black] (0.5,0.4625)--(2.5,1.9625);
 \draw[thick,blue!70!black,domain=0.15:2.65,samples=80] plot (\x,{0.4+0.25*\x*\x});
 \fill (0.5,0.4625) circle (1.5pt);
 \fill (2.5,1.9625) circle (1.5pt);
 \node[below] at (0.5,0) {$a$};
 \node[below] at (2.5,0) {$b$};
 \node[above left] at (1.55,1.25) {$L(x)$};
 \node[right] at (2.05,1.18) {$f(x)$};
\end{tikzpicture}
\caption{端点を結ぶ直線の積分が台形公式になる。下に凸の曲線では，弦と曲線の間の面積だけ過大になる。}
\label{図:台形近似}
\end{figure}

$w(x)=(x-a)(b-x)$と置く。端点で$w$が$0$となり，$w''(x)=-2$となることを利用して，二回部分積分する。
\[
\begin{aligned}
 \int_a^b w(x)f''(x)\dd x
 &=-\int_a^bw'(x)f'(x)\dd x\\
 &=-[w'(x)f(x)]_a^b+\int_a^bw''(x)f(x)\dd x\\
 &=(b-a)\{f(a)+f(b)\}-2\int_a^bf(x)\dd x.
\end{aligned}
\]
したがって，誤差を「積分の値から近似値を引いたもの」と決めれば
\[
 \boxed{\int_a^bf(x)\dd x-T
 =-\frac12\int_a^b(x-a)(b-x)f''(x)\dd x}.
\]
$f''\geq0$なら右辺は$0$以下になるので，図で見た大小関係が式からも分かる。また$|f''(x)|\leq M_2$なら
\[
 \left|\int_a^bf(x)\dd x-T\right|
 \leq\frac{M_2}{2}\int_a^b(x-a)(b-x)\dd x
 =\frac{M_2(b-a)^3}{12}.
\]

区間全体を$N$等分し，$h=\frac{b-a}{N}$，$x_j=a+jh$とすると，各小区間の台形を足した値は
\[
 T_N=h\left\{\frac{f(a)+f(b)}2+\sum_{j=1}^{N-1}f(x_j)\right\}.
\]
内部の点は隣り合う二つの台形で半分ずつ使われるため，係数が$1$になる。各小区間の誤差は$\frac{M_2h^3}{12}$以下なので，合計では
\[
 \left|\int_a^bf(x)\dd x-T_N\right|
 \leq\frac{M_2(b-a)h^2}{12}
 =\frac{M_2(b-a)^3}{12N^2}.
\]
分割数を二倍にすれば，この誤差上限は四分の一になる。

